
%% \begin{figure*}[t]
%%   \label{fig-cx-star}
%%   \begin{center}
%%   \end{center}
%%   \caption{The $C-X^{*}$ gate from \cite{bib-amy-phase-state}}
%% \end{figure*}
%% \begin{figure*}[t]
%%   \begin{minipage}{.46\linewidth}[t]
%%     \label{fig-lut-network}
%%     \begin{center}
%%       \Qcircuit @C=0.5em @R=0.2em @!R { \\
  \lstick{\ket{{x}_{1}}} & \qw  & \qw                 & \qw                 & \multigate{5}{U_3}      & \qw                & \qw                   & \qw                          & \qw                          & \qw                           & \ghost{x_1}  & \lstick{\ket{{x}_{1}}} \\
  \lstick{\ket{{x}_{2}}} & \qw  & \multigate{1}{U_1}  & \qw                 & \push{\dashedwire}\qw   & \qw                & \qw                   & \qw                          & \qw                          & \multigate{1}{U_1^{\dagger}}  & \ghost{x_2}  & \lstick{\ket{{x}_{2}}} \\
  \lstick{\ket{{x}_{3}}} & \qw  & \ghost{1}           & \qw                 & \push{\dashedwire}\qw   & \qw                & \qw                   & \qw                          & \qw                          & \ghost{U_1^{\dagger}}         & \ghost{x_3}  & \lstick{\ket{{x}_{3}}} \\
  \lstick{\ket{{x}_{4}}} & \qw  & \qw                 & \multigate{1}{U_2}  & \push{\dashedwire}\qw   & \qw                & \qw                   & \qw                          & \multigate{1}{U_2^{\dagger}} & \qw                           & \ghost{x_4}  & \lstick{\ket{{x}_{4}}} \\
  \lstick{\ket{{x}_{5}}} & \qw  & \qw                 & \ghost{2}           & \push{\dashedwire}\qw   & \qw                & \multigate{4}{U_5}    & \qw                          & \ghost{U_2^{\dagger}}        & \qw                           & \ghost{x_5}  & \lstick{\ket{{x}_{5}}} \\
  \lstick{\ket{0}     }  & \qw  & \targ\qwx[-3]\qw    & \qw                 & \ghost{3}\qw            & \multigate{1}{U_4} & \push{\dashedwire}\qw & \multigate{1}{U_4^{\dagger}} & \qw                          & \targ\qwx[-3]\qw              & \ghost{0}    & \lstick{\ket{0}     }    \\
  \lstick{\ket{0}     }  & \qw  & \qw                 & \targ\qwx[-2]\qw    & \qw                     & \ghost{4}          & \push{\dashedwire}\qw & \ghost{U_4^{\dagger}}        & \targ\qwx[-2]\qw             & \qw                           &  \qw            & \lstick{\ket{0}     }  \\
  \lstick{\ket{0}     }  & \qw  & \qw                 & \qw                 & \targ\qwx[-2]\qw        & \qw                & \push{\dashedwire}\qw & \qw                          & \qw                          & \qw                           & \ghost{y_1}  & \lstick{\ket{y_1}   }  \\      
  \lstick{\ket{0}     }  & \qw  & \qw                 & \qw                 & \qw                     & \targ\qwx[-2]\qw   & \ghost{5}             & \targ\qwx[-2]\qw             & \qw                          & \qw                           & \ghost{0}    & \lstick{\ket{0}     }    \\
  \lstick{\ket{0}     }  & \qw  & \qw                 & \qw                 & \qw                     & \qw                & \targ\qwx[-1]\qw      & \qw                          & \qw                          & \qw                           & \ghost{y_2}  & \lstick{\ket{y_2}     }  \\
}

%%     \end{center}
%%   \end{minipage}
  
%% \end{figure*}
%\vspace*{0.25cm}

\section{Proposed Method}
\label{Sec:proposed}
\subsection{Realizing Arbitrarily Controlled NOT Functions Up To A Relative Phase}
\label{Subsec:UNOT}

\begin{figure}[tbp]
    \centering
    \scalebox{0.8}{\Qcircuit @C=0.5em @R=0.2em @!R { \\
          \push{{x}_{0} }   & & \multigate{3}{x_k \cdot U}  & \qw   & \qw & &         & & \qw & \qw               & \qw                       & \qw      & \qw                       & \multigate{3}{U}              & \qw               & \qw      & \qw                         & \qw               & \qw & \\
          \push{{x}_{1} }   & & \ghost{x_k \cdot f}         & \qw   & \qw & &         & & \qw & \qw               & \qw                       & \qw      & \qw                       & \ghost{U}                      & \qw               & \qw      & \qw                         & \qw               & \qw & \\
          \push{\vdots}     & &                             &       &     & &\push{=} & &     &                   & \push{\vdots}             &          &  \push{\vdots}            &                                & \push{\vdots}     &          &                             & \push{\vdots}     &     & \\
          \push{{x}_{n-1} } & & \ghost{x_k \cdot f}         & \qw   & \qw & &         & & \qw & \qw               & \qw                       & \ctrl{1} & \qw                       & \ghost{U}                      & \qw               & \ctrl{1} & \qw                         & \qw               & \qw & \\
          \push{{x}_{n} }   & & \push{\rtoftarg}\qw\qwx[-1] & \qw   & \qw & &         & & \qw & \gate{\mathrm{H}} & \gate{\mathrm{T}}         & \targ    & \gate{\mathrm{T^\dagger}} & \targ\qwx[-1]    & \gate{\mathrm{T}} & \targ    & \gate{\mathrm{T^{\dagger}}} & \gate{\mathrm{H}} & \qw & \qw\\
}
                            
}
    \caption{A U-controlled RTOF gate}
    \label{fig-urtof}
\end{figure}

To get to a point where we can realize arbitrary Boolean functions up to a relative phase, we first observe that the RTOF depicted in Fig.~\ref{fig-rtof} can replace the center CNOT with a $n$-input Toffoli gate or U-controlled NOT (UCNOT) as depicted in Fig.~\ref{fig-urtof} which was observed by Maslov~\cite{bib-rtof-maslov}. However, as noted in~\cite{bib-amy-phase-state}, the introduction of a relative phase here has an effect on the final Boolean function as the final Hadamard gate rotates the relative phase back into state space. Therefore, we take the $C-X^*$ construction that is depicted in~\cite{bib-amy-phase-state}. We can generalize this $n$-input Toffoli gate to an arbitrarily controlled NOT gate as depicted in Fig.~\ref{fig-urtof}. 

\begin{figure}[tb]
    \centering
    \scalebox{0.8} {
    \Qcircuit @C=0.5em @R=0.2em @!R { \\
          \push{{x}_{0} }     & & \multigate{3}{D} & \multigate{3}{x_k \cdot g} & \qw   & \qw & &         & & \qw               & \qw                     & \multigate{3}{g} & \qw                     & \multigate{3}{f}  & \qw               & \multigate{3}{g} & \qw                        & \multigate{3}{f}   & \qw \\
          \push{{x}_{1} }     & & \ghost{D}        & \ghost{x_k \cdot g}        & \qw   & \qw & &         & & \qw               & \qw                     & \ghost{g}\qw     & \qw                     & \ghost{f}         & \qw               & \ghost{g}        & \qw                        & \ghost{f}          & \qw \\
          \push{\vdots}      & &                  &                            &       & \qw & &\push{=} & &                    & \push{\vdots}           &                  &  \push{\vdots}         &                   & \push{\vdots}     &                  &                            &                   & \push{\vdots} \\
          \push{{x}_{n} }     & & \ghost{D}        & \ghost{x_k \cdot g}       & \qw   & \qw & &         & & \qw               & \qw                     & \ghost{g}\qw     & \qw                     & \ghost{f}         & \qw               & \ghost{g}        & \qw                        & \ghost{f}          & \qw \\
          \push{{\ket{0}}   }& & \ghost{D}        & \targ\qwx[-1]              & \qw   & \qw & &         & & \gate{\mathrm{H}} & \gate{\mathrm{T^\dagger}} & \targ\qwx[-1]    & \gate{\mathrm{T^\dagger}} & \targ\qwx[-1]    & \gate{\mathrm{T}} & \targ\qwx[-1]    & \gate{\mathrm{T^{\dagger}}} & \qw                 &\gate{\mathrm{H}} \\
}


    }
    \caption{Generalization of Fig.~\ref{fig-urtof}, which also implements the Boolean product of two functions}
    \label{fig-f-dot-g}
\end{figure}

We can further generalize this. Recall from Section~\ref{Sec:Exist} that a Toffoli gate realizes a Boolean product of its inputs on its output. Just like above, we can use the same reasoning to generalize both inputs into arbitrary Boolean functions~\cite{bib-amy-phase-state} . This means that we can use the construction in Fig.~\ref{fig-f-dot-g} to realize a product between two arbitrary Boolean functions up to a relative phase. 
% \begin{figure*}[t]
%   \begin{minipage}{0.45\linewidth}
%     \begin{center}
%       \scalebox{0.5} {
%         \Qcircuit @C=0.5em @R=0.2em @!R { \\
          \push{{x}_{0} }     & & \multigate{5}{U_F}             & \qw  & \qw & &         & & \qw               & \qw               & \qw           & \qw                       & \multigate{5}{U_{F_{x_k=1}}}    & \qw               & \qw           & \qw                         & \multigate{5}{U_{F_{x_k=1}}^{\dagger}} & \qw           & \qw & \qw               & \qw             & \qw                       & \multigate{5}{U_{F_{x_k=0}}} & \qw               & \qw             & \qw                          & \multigate{5}{U_{F_{x_k=0}^{\dagger}}}    & \qw               & \qw\\
          \push{{x}_{1} }     & & \ghost{U_F}                    & \qw  & \qw & &         & & \qw               & \qw               & \qw           & \qw                       & \ghost{U_{F_{x_k=1}}}\qw        & \qw               & \qw           & \qw                         & \ghost{U_{F_{x_k=1}}^{\dagger}}        & \qw           & \qw & \qw               & \qw             & \qw                       & \ghost{U_{F_{x_k=0}}}        & \qw               & \qw             & \qw                          & \ghost{U_{F_{x_k=0}}^{\dagger}}\qw        & \qw               & \qw\\
          \push{\vdots}       & &                                &      & \qw & &         & &                   & \push{\vdots}     &               &  \push{\vdots}            &                                 & \push{\vdots}     &               &                             &                                        & \push{\vdots} &     &                   &                 &  \push{\vdots}            &                              & \push{\vdots}     &                 &                              &                                           & \push{\vdots}     &    \\
          \push{{x}_{k} }     & & \ghost{U_F}                    & \qw  & \qw & &         & & \qw               & \qw               & \ctrl{3}      & \qw                       & \ghost{U_{F_{x_k=1}}}\qw        & \qw               & \ctrl{3}      & \qw                         & \ghost{U_{F_{x_k=1}}^{\dagger}}        & \qw           & \qw & \gate{X}          & \ctrl{3}        & \qw                       & \ghost{U_{F_{x_k=0}}}\qw     & \qw               & \ctrl{3}        & \qw                          & \ghost{U_{F_{x_k=0}}^{\dagger}}\qw        & \gate{X}          & \qw\\
          \push{\vdots}       & &                                &      & \qw & &\push{=} & &                   & \push{\vdots}     &               & \push{\vdots}             &                                 & \push{\vdots}     &               &                             &                                        & \push{\vdots} &     &                   &                 &  \push{\vdots}            &                              & \push{\vdots}     &                 &                              &                                           & \push{\vdots}     &    \\
          \push{{x}_{n} }     & & \ghost{U_F}                    & \qw  & \qw & &         & & \qw               & \qw               & \qw           & \qw                       & \ghost{U_{F_{x_k=1}}}\qw        & \qw               & \qw           & \qw                         & \ghost{U_{F_{x_k=1}}^{\dagger}}        & \qw           & \qw & \qw               & \qw             & \qw                       & \ghost{U_{F_{x_k=0}}}\qw     & \qw               & \qw             & \qw                          & \ghost{U_{F_{x_k=0}}^{\dagger}}\qw        & \qw               & \qw\\
          \push{{\ket{0}}   } & & \push{\rtoftarg}\qw\qwx[-1]    & \qw  & \qw & &         & & \gate{\mathrm{H}} & \gate{\mathrm{T}} & \targ\qwx[-3] & \gate{\mathrm{T^\dagger}} & \push{\rtoftarg}\qw\qwx[-1]     & \gate{\mathrm{T}} & \targ\qwx[-3] & \gate{\mathrm{T^{\dagger}}} & \push{\rtoftarg}\qw\qwx[-1]            & \qw           & \qw & \gate{\mathrm{T}} & \targ\qwx[-1]   & \gate{\mathrm{T^\dagger}} & \push{\rtoftarg}\qw\qwx[-1]  & \gate{\mathrm{T}} & \targ\qwx[-1]   & \gate{\mathrm{T^{\dagger}}}  & \push{\rtoftarg}\qw\qwx[-1]               & \gate{\mathrm{H}} & \qw\\
}
%                              & &                                &      &     & &         & & \push{$g_0$}      & \push{$g_1$}      & \push{$g_2$}  & \push{$g_3$}              & \push{$g_4$}                & \push{$g_5$}      & \push{$g_6$}  & \push{$g_7$}                & \push{$g_8$}                 &               &     & \push{$g_9$}      & \push{$g_{10}$} & \push{$g_{11}$}           & \push{$g_{12}$}             & \push{$g_{13}$}   & \push{$g_{14}$} & \push{$g_{15}$}                     & \push{$g_{16}$}                    & \push{$g_{17}$}

%       }
%     \end{center}
%     \caption{Balanced Shannon's Decomposition}
%     \label{fig-shannon}
%   \end{minipage}
%   \hspace{0.5cm}
%   \begin{minipage}{0.45\linewidth}
%     \begin{center}
%       \scalebox{0.5} {
%         \Qcircuit @C=0.5em @R=0.2em @!R { \\
          \push{{x}_{0} }     & & \multigate{5}{F} & \qw  & \qw & &         & & \multigate{5}{F_{x_k=0}} & \qw               & \qw               & \qw            & \qw                       & \multigate{5}{F_{x_k=0} \oplus F_{x_k=1}} & \qw               & \qw           & \qw                          & \qw               & \qw\\
          \push{{x}_{1} }     & & \ghost{F}        & \qw  & \qw & &         & & \ghost{F_{x_k=0}}\qw     & \qw               & \qw               & \qw            & \qw                       & \ghost{F_{x_k=0} \oplus F_{x_k=1}}        & \qw               & \qw           & \qw                          & \qw               & \qw\\
          \push{\vdots}       & &                  &      & \qw & &         & &                          &                   &                   &                &  \push{\vdots}            &                                           & \push{\vdots}     &               &                              & \push{\vdots}     &    \\
          \push{{x}_{k} }     & & \ghost{F}        & \qw  & \qw & &         & & \ghost{F_{x_k=0}}\qw     & \qw               &                   & \ctrl{3}       & \qw                       & \ghost{F_{x_k=0} \oplus F_{x_k=1}}\qw     & \qw               & \ctrl{3}      & \qw                          &                   & \qw\\
          \push{\vdots}       & &                  &      & \qw & &\push{=} & &                          &                   &                   &                &  \push{\vdots}            &                                           & \push{\vdots}     &               &                              & \push{\vdots}     &    \\
          \push{{x}_{n} }     & & \ghost{F}        & \qw  & \qw & &         & & \ghost{F_{x_k=0}}\qw     & \qw               & \qw               & \qw            & \qw                       & \ghost{F_{x_k=0} \oplus F_{x_k=1}}\qw     & \qw               & \qw           & \qw                          & \qw               & \qw\\
          \push{{\ket{0}}   } & & \targ\qwx[-1]    & \qw  & \qw & &         & & \targ\qwx[-1]            & \gate{\mathrm{H}} & \gate{\mathrm{T}} & \targ\qwx[-1]  & \gate{\mathrm{T^\dagger}} & \targ\qwx[-1]                             & \gate{\mathrm{T}} & \targ\qwx[-1] & \gate{\mathrm{T^{\dagger}}}  & \gate{\mathrm{H}} & \qw\\
}

%       }
%     \end{center}
%     \caption{Unbalanced Positive Davio Decomposition}
%     \label{fig-davio-pos-unb}
%   \end{minipage}
% %\vspace*{0.5cm}
% \end{figure*}

\subsection{Using Shannon's Decomposition To Realize Arbitrarily Controlled Single Target Gates}
\label{Subsec:shannon}

\begin{figure}[tbp]
    \centering
      \scalebox{0.7} {
        \Qcircuit @C=0.5em @R=0.2em @!R { \\
          \push{{x}_{0} }  & \multigate{5}{x_k \cdot F} & \qw      & \multigate{5}{\bar{x_k} \cdot F'}&\qw\\
          \push{{x}_{1} }  & \ghost{x_k \cdot F}        & \qw      & \ghost{\bar{x_k} \cdot F'}       &\qw\\
          \push{\vdots}    &                          &          &                                &   \\
          \push{{x}_{k} }  & \ghost{x_k \cdot F}        & \gate{X} & \ghost{\bar{x_k} \cdot F'}       &\qw\\
          \push{\vdots}    &                          &          &                                &   \\
          \push{{x}_{n} }  & \ghost{x_k \cdot F}        & \qw      & \ghost{\bar{x_k} \cdot F'}       &\qw\\
          \push{{\ket{0}}} & \targ\qwx[-1]            & \qw      & \targ\qwx[-1]                  &\qw\\
}

      }
    \caption{A generalized way of conducting the OR for the Shannon decomposition}
    \label{fig-bool-or}
\end{figure}

Recall Shannon's decomposition $f(x_n,...,x_1) = x_n \cdot f(x_n=1,x_{n-1},\cdots,x_0) + \bar{x_n} \cdot f(x_n=0,x_{n-1},\cdots,x_0)$.
Using the constructions mentioned in Section~\ref{Subsec:UNOT}, we thus have a way to recursively construct
arbitrary Boolean functions up to a relative phase as a single target controlled NOT gate.

As a prelude to our discussion, we observe that $x \cdot a  + \bar{x} \cdot b = x \cdot a \oplus \bar{x} \cdot b$.
This means that we can implement this function as a series of Toffoli gates. Generalizing to UCNOT gates, we get
the construction shown in Fig.~\ref{fig-bool-or}, that we will be using to construct our implementations in the
next two sections.


First, from the Shannon's decomposition
formula~$f(x_n,...,x_1) = x_n \cdot f(x_n=1,x_{n-1},\cdots,x_0) + \bar{x_n} \cdot f(x_n=0,x_{n-1},\cdots,x_0)$
where we can take $g=f(x_n=1,x_{n-1},\cdots,x_0)$ and $g^{\prime}=f(x_n=0,x_{n-1},\cdots,x_0)$. As always with
Shannon's decomposition, we can decompose using an arbitrary variable order. (We will defer discussion of how we
determine this variable order to Section~\ref{Sec:exp}). We take Fig.~\ref{fig-f-dot-g} to insert our function, and
we get the construction as shown in Fig.~\ref{app-shannon}. We then take Fig.~\ref{fig-bool-or} and replace each
UCNOT in it with one of the terms in the Shannon's decomposition that we just made. As noted in
Section~\ref{Subsec:UNOT}, it is easy to see that this expression implements the Sum of Products (SOP) expression
of Shannon's decomposition. We then take each of $g$ and $g^{\prime}$ and repeat the process recursively until we
are left with a simple CNOT as our U-controlled NOT expression.

If it is ``False'', we have a False-controlled-NOT, meaning an identity function. This means that if one side of
the Shannon decomposition evaluates to either ``True'' or ``False'', we can merely replace it with a NOT gate or an
identity (meaning nothing at all), respectively, without expending any additional T-gates to implement the
function. As we will see in the next example, we can express the relative-phase Toffoli gate construction
in~\cite{bib-amy-phase-state} as a special case of this construction

Note that with this decomposition, when the function evaluates as ``True'', we end up with a True-controlled-NOT,
meaning, a NOT function. Because the Shannon decomposition now evaluates to either $x_k \cdot F_{x=1} + \bar{x}$ or
$x_k + \bar{x_k} \cdot F_{x=0}$, we can just implement the side with a True expression using an $x_k$ or
$\bar{x_k}$ controlled~CNOT gate. 

\begin{example}
  \label{ex-bool}
  We attempt to implement the function $F(x_3,x_2,x_1) = x_3x_2x_1 + \bar{x_3}(x_1 + x_2)$ as its relative phase equivalent.

  For the $F(x_3 = 1,x_2,x_1)$ case, we note that this construction just looks like a 3-input Toffoli gate. The first stage of this decomposition looks just like the
  decomposition in Fig~\ref{rtof}. We can see that if we
  continue Shannon's decomposition, one or the other side of the
  equation will always evaluate to 0. Therefore it is easy to see that it simplifies to the
  decomposition in Fig~\ref{fig-f-dot-g-in}, which is the decomposition in ~\ref{bib-amy-phase-state}.
  (As an aside, it generalizes it slightly as the algorithm also automatically handles the negated input
  case without any additional information). %The result is labeled with a curly brace in Fig.~\ref{fig-func-decomp}. 
  
  For the $F(x_3 = 1,x_2,x_1)$ case, we now have something that is not a straightforward sum of products term.
  We now move on to the next level (Note: The variable order here, $x_3 \rightarrow x_2 \rightarrow x_1$, is
  chosen to best illustrate the working of the method. We treat optimal variable order later). We see that
  $F(x_3 = 0,x_2=0,x_1) = x_1$. Following the algorithm here results in something that looks like an $x_2$
  and $x_1$ controlled Toffoli gate. We then move on to $F(x_3 = 0,x_2=1,x_1)$ and find that it is $True$.
  This is one of the special cases in our algorithm, and so we can implement $x_2 \cdot F(x_3 = 0,x_2=1,x_1)$
  as an $x_2$ controlled NOT gate. We then use the construction depicted in Fig~\ref{fig-bool-or} to compose
  the two functions. %The result is depicted Fig.~\ref{fig-func-decomp}, indicated by the dashed box. 
  
  Finally we compose $F(x_3 = 1,x_2,x_1)$ and $F(x_3 = 1,x_2,x_1)$ as in Fig~\ref{fig-bool-or} again. The result is depicted in Fig.~\ref{app-func-decomp} in ~\ref{App:FuncDecomp}.
\end{example}

However, as we will see later in Section \ref{Sec:Exp}, some expressions are not amenable to a purely Shannon
Decomposition. Several of them have highly complex (many BDD nodes) representation in the Shannon Decomposition.
We must then introduce a different representation to take them into account.

\subsection{Integrating Davio Expansions}
\label{Subsec:davio}
As can be seen in the previous section, a significant portion of the cost comes from multiplication.
Therefore it is useful to consider a method that takes advantage of {\it Exclusive Sum of Products} (ESOP)
expressions. Recall the positive Davio Expansion $F=F_1\oplus x_n \cdot F_2$ and the corresponding
negative Davio expansion $F=F_0 \oplus \bar{x_n} \cdot F_2$ where $F_0 = F(x_n=0)$,$F_1 = F(x_n=1)$,
and $F_2 = F_0 \oplus F_1$. These expansions have the advantage that at every level, only one
multiplication takes place, which already reduces the maximum possible T-gate use at any node from
8 to 4. 

\subsection{Using Balanced and Unbalanced Relative Phase Functions}
\label{Subsec:balanced}
As Amy et al. ~\cite{bib-amy-phase-state} pointed out, because any multiplication is first ``constructed''
in phase space and then rotated back into state space with a Hadamard gate, adding a relative phase circuit
inside that logic with a \emph{target dependent} phase (meaning one whose value depends on the target bit
of the function) in an \emph{unbalanced} construction, such as in Fig~\ref{fig-urtof}, will see its relative
phase added on to the main Boolean function as a residue i.e. it will take the form of $F=x \cdot F_k + F_{k-1}$.
To keep this from happening, a \emph{balanced} construction such as in Fig~\ref{fig-f-dot-g} is used, with pairs
of relative phase circuits and their inverses to cancel out the target dependence that the Hadamard gate on the
bit will act on. In the process, such a construction creates a function with a \emph{target independent}relative
phase. This means that it can be used inside an unbalanced construction without incurring residue since the
Hadamard gate now does not have any dependence to act on.

Amy et al. propose a method where balanced and unbalanced constructions are alternated at different levels of the
construction of an MCT to reduce the amount of times a function has to be doubled to cancel itself out. We can
easily extend this to our method of constructing arbitrary Boolean functions above. We start at the root node
with an unbalanced construction, and any children that are constructed recursively will have to be constructed
in a balanced manner. The children of those in turn can be unbalanced, and so on. 

In doing so, we can reduce the dependency from $2^n$ to $2^{\frac{n}{2}}$. While still exponential, this has quite
a significant reduction in the regime we are working with. For example for $n=6$, the circuit is at MOST 1/8th of
the original cost! Merely introducing this methodology results in significant savings for larger LUT-sizes, though
its effect could be diminished for smaller ones.

This has additional ramifications when combined with the Davio Expansions. Since there is only one multiplication
performed, the non-multiplied term ($F_0$ in the positive Davio case and $F_1$ in the negative Davio case) can
be implemented as an \emph{unbalanced} relative phase circuit, provided the parent construction can also be
unbalanced. This means there is no need to double the logic inside to cancel the residue, and so the exponential
is further reduced. This can reduce the exponential relation even further.

\subsection{Relative Phase Synthesis Algorithm}
\label{Subsec:Pro:Improved}
We now try to integrate the Sec~\ref{Subsec:shannon}, Sec~\ref{Subsec:balanced}, and ~\ref{Subsec:davio}.

As a prelude to our discussion, we observe that $x \cdot a  + \bar{x} \cdot b = x \cdot a \oplus \bar{x} \cdot b$.
This means that we can implement this function as a series of Toffoli gates. Generalizing to UCNOT gates, we get
the construction shown in Fig.~\ref{fig-bool-or}.

First, we determine the variable order to decompose the function in. As the details of the variable order
computation are not important to the algorithm itself, we will defer discussion of how we determine
this variable order to our discussion of the implementation in Section~\ref{Sec:exp}.

If the function is neither ``True'' nor ``False'', the function is expanded using the Shannon, Positive
Davio, or Negative Davio expansions, independently. We go over constructions for them in Appendix
~\ref{App:ManyDecomp}. The constituent subfunctions of each expansion are then processed in the same way by
recursing the function with the decomposed Boolean function as input.. When the logic can be implemented as an
unbalanced function, such as in the unmultiplied term of the Davio expansion or as balanced logic gates, that is
relayed to the next level of the recursion via an input.

If it is ``False'', we have a False-controlled-NOT, meaning an identity function. In this case we remove the term
with this function entirely.

Note that with this decomposition, when the function evaluates as ``True'', we end up with a True-controlled-NOT, meaning, a NOT function. With the multiplied terms such as the $x_k \cdot F_{x=1}$ in Shannon's Decomposition or $\bar{x_k} \cdot (F_{x=0} \oplus F_{x=1})$ in the Davio Expansion, we simply get $x_k$ or $\bar{x_k}$ respectively. This means we can just implement the side with a True expression using an $x_k$ or $\bar{x_k}$ controlled~CNOT gate. 

As the recursion terminates, the cost of each of the Shannon, Positive Davio, or Negative Davio expansions is
compared with each other, and the minimum is chosen. The result is then passed on to the recursion level above it.
The process is then repeated with all of the variable orders and the minimum is picked from the results. The
resulting algorithm is described in Alg.~\ref{algo-new}


\begin{algorithm}[tb]
  \caption{relative\_phase(BF: a Boolean function, balanced: Boolean indicating whether this function is balanced)}
  \label{algo-new}
  \begin{algorithmic}
%    \State BF $\leftarrow$ A Boolean Function
    \State QC $\leftarrow$ an empty quantum circuit with registers
    \State Output $\leftarrow$ the output bit of the Boolean Function
    \State order $\leftarrow$ BF.find\_optimal\_order()
    \If{$BF == True$}
    \State return CNOT
    \ElsIf{$BF == False$}
    \State return QC $\leftarrow$ Identity(Output)
    \ElsIf{$BF$ is a leaf node}
    \State return CNOT
    \Else
    \State    //The Shannon $F_0$ and $F_1$ alternate balanced and unbalanced
    \State    $F_0$ $\leftarrow$ relative\_phase($BF_{x_{order[first]}\leftarrow 0)}$,!balanced) 
    \State    $F_1$ $\leftarrow$ relative\_phase($BF_{x_{order[first]}\leftarrow 1}$,!balanced) 
    \State
    \State    //The Davio $F_0$ and $F_1$ inherit the parent's balanced attribute
    \State    $(F_0)_{dav}$ $\leftarrow$ relative\_phase($BF_{x_{order[first]}\leftarrow 0)}$,balanced) 
    \State    $(F_1)_{dav}$ $\leftarrow$ relative\_phase($BF_{x_{order[first]}\leftarrow 1}$,balanced) 
    \State
    \State    $F_2$ $\leftarrow$ relative\_phase($BF_{x_{order[first]}\leftarrow 1} \oplus BF_{x_{order[first]}\leftarrow 0}$,
    \State                       !balanced)
    \State
    \State    $shannon$ $\leftarrow$ compose $F_0$ and $F_1$ in a balanced manner if balanced or unbalanced if not.
    \State    $davio_{pos}$ $\leftarrow$ compose $(F_0)_{dav}$ and $F_2$ in a balanced manner if balanced is true or unbalanced if not.
    \State    $davio_{neg}$ $\leftarrow$ compose $(F_1)_{dav}$ and $F_2$ in a balanced manner if balanced is true or unbalanced if not.
    \State    return one of $[shannon, davio_{pos},davio_{neg}]$ with the min T-count
    \EndIf
    \State Repeat for every variable order and pick the minimum
  \end{algorithmic}
  
\end{algorithm}


\begin{example}
  \label{ex-davio}
    We again attempt to implement the function $F(x_3,x_2,x_1) = x_3x_2x_1 + \bar{x_3}(x_1 + x_2)$ as its relative phase equivalent.
    Just as before, we first take the positive Davio expansion about $x_3$ (for simplicity, we will only follow the Davio branches
  in this example. However, as  we will see later, they turn out to be optimal for this function). We get $F_0 = x_1 + x_2$ and
  $F_2 = F_0 \oplus F_1 = (x_1 + x_2) \oplus x_1 \cdot x_2 = x_1 \oplus x_2$. Note that for this level we can tolerate an unbalanced
  function, so for the $F_0$ branch we set the ``balanced'' parameter to false and on the $F_2$ branch we set it to true.

  We then proceed to $x_2$. We decompose $F_0 = (F_0)_0 \oplus x_2 \cdot ((F_0)_2)$. Since $(F_0)_0 = x_1$, we just implement it as an $x_1$
  controlled CNOT. After calculating $(F_0)_1 = x_1 + 1 = 1$, we get $(F_0)_2 = (F_0)_0 \oplus (F_0)_1 = x_1 \oplus 1 = \bar{x_1}$
  Thus $x_2 \cdot (F_0)_2$ is implemented according to a balanced function multiplication such as in Fig.~\ref{fig-f-dot-g}. We then compose
  these similarly to Fig.~\ref{app-davio-pos-unb} but with a balanced constructionand terminate the recursion. We can see at this point that
  there are 4 T-gates.

  We now go to the $F_2$ branch. Here, since $x_1 \oplus x_2$ can be implemented using an $x_1$ controlled CNOT and $x_2$ controlled CNOT
  in series, we synthesize that and terminate the recursion.

  Back at the top level, we compose the subcircuits for $F_0$ and $F_2$ as a negative unbalanced Davio construction, adding an additional 4
  T-gates. This brings us to a total of 8 T-gates. Compare this to the 24 T-gates of our original methodology, 1/3 the size. We omit
  the computation of the other branches and variable orders in the interest of space, but the reader can perform them  quickly to
  see that the branch we followed is the optimal branch (or one of them. In this function there are multiple because $x_1$  and
  $x_2$ are symmetric, and several Davio branches have similar costs). 
  The resultant circuit is depicted in Fig.~\ref{fig-func-decomp-davio} in \ref{App:FuncDecomp}

\end{example}

\begin{example}

%% \begin{algorithm}[tbp]
%%   \caption{shannon\_decomp(BF: a Boolean function)}
%%   \label{algo-shannon}
%%   \begin{algorithmic}
%% %    \State BF $\leftarrow$ A Boolean Function
%%     \State QC $\leftarrow$ an empty quantum circuit with registers
%%     \State Output $\leftarrow$ the output bit of the Boolean Function
%%     \State order $\leftarrow$ BF.find\_optimal\_order()
%%     \If{$BF == True$}
%%     \State return CNOT
%%     \ElsIf{$BF == False$}
%%     \State return QC $\leftarrow$ Identity(Output)
%%     \ElsIf{$BF$ is a leaf node}
%%     \State return CNOT
%%     \Else
%%     \State $F_0$ $\leftarrow$ shannon\_decomp($BF_{x_{order[first]}\leftarrow 0)}$)
%%     \State $F_1$ $\leftarrow$ shannon\_decomp($BF_{x_{order[first]}\leftarrow 1}$)
%%     \State return a UCNOT implemented by composing $F_0$ + $F_1$ in the manner shown in Fig.~\ref{app-shannon}
%%     \EndIf
%%   \end{algorithmic}
  
%% \end{algorithm}

\begin{algorithm}[tb]
  \caption{lut\_synthesis}
  \label{algo-lut}
  \begin{algorithmic}
    \State lut\_net $\leftarrow$ a LUT network
    \State esop (Boolean Function) $\leftarrow$ a method to generate
    an exact QC to calculate a Boolean function implemented as a UCNOT gate 
    \For{every node in lut\_net}
    \If{node is intermediate node}
    \State QC $\leftarrow$ add relative\_phase(node's Boolean function) on a dirty ancilla
    \State QC $\leftarrow$ add inverse of relative\_phase(node's 
    Boolean function) where we restore the state of the same dirty ancilla
    \ElsIf{node is output node}
    \State QC $\leftarrow$ add esop (node's Boolean function) acting on a dirty ancilla
    \EndIf
    \EndFor
  \end{algorithmic}
%\vspace*{-0.25cm}
\end{algorithm}

\subsection{Synthesizing A Boolean Function Using LUT-networks and Shannon Decomposition RTOF}

Using Algorithm~\ref{algo-new} we are now ready to propose our LUT-based synthesis method 

First we compose the circuit as in Sec.~\ref{Sec:lhrs}. However, this time, we replace the alternating pairs of compute/uncompute gates with the respective relative phase version using Algorithm~\ref{algo-shannon} as well as its associated inverse. If there is no matching uncompute logic i.e., a node acting on an output qubit, we generate the logic exactly. The method is enumerated explicitly in Algorithm~\ref{algo-lut}.

This QC is then parsed by our algorithm, synthesizing each UCNOT into relative phase versions of the partial Boolean functions in each intermediate node (as well as their inverse), and exact versions in each output node. These are then recombined into a quantum circuit. 

\label{ex-lut-rp}
  We follow the steps in Ex.~\ref{ex-lut}. We then parse through LUTs
  $1$,$2$,$1^{\dagger}$,$2^{\dagger}$,$4$,and $4^{\dagger}$ and
  generate their relative phase versions. For nodes 3 and 5, we use an
  ESOP-based exact synthesis. 
  The result is depicted in Fig.~\ref{fig-lut-qc}.
\end{example}


